%قبل از شروع مسئله ترجیحاً چیزی ننویسید. اما اگر خواستید بنویسید، در قالب لیستی (فقرات و شمارش) بنویسید.
\شروع‌مسئله[20(امتیازی)]  
% صورت اصلی مسئله را در این‌جا بنویسید.
\\
یک جدول درهم سازی از اندازه‌ی $m$  در نظر بگیرید که $n$ آیتم را در خودش ذخیره کرده است که$n \le m/2$. برای رفع مشکل برخود نیز از روش آدرس دهی باز استفاده شده است. فرض کنید متغیر تصادفی $ x_i$  نمایانگر تعداد جستجو های موردنیاز  برای $i$  امین درج کردن باشد. فرض کنید متغیر تصادفی $X$  برابر بیشینه مقدار $x_i$ ها باشد. ثابت کنید که $E[X]$  از مرتبه‌ی $ O(log(n))$ است. (راهنمایی :‌ ثابت کنید $probability \{X > 2log(n)\} \le \frac{1}{n}$ سپس از این موضوع در محاسبه‌ی $E[X]$ و تخمین های مربوطه استفاده کنید .) 


\solution{
اگر متغیر تصادفی X نمایانگر تعداد جستجو های برای یافتن کلید در یک سری جستجوی ناموفق باشد آنگاه\\
$Pr\{X \ge i\} \le \alpha^{i-1}$ (چرا؟). 
چون $n \le m/2$ پس داریم: $\alpha \le 1/2$. حال $i$ را برابر $k+1$ قرار می‌دهیم پس داریم:
\\
\begin{LTR}
$Pr\{X > k\} = Pr\{X \ge k+1\} \le (1/2)^{(k+1)-1} = 2^{-k} .$
\end{LTR}{}

\\
حال 
$k$
را برابر
$2logn$
قرار می‌دهیم، پس احتمال اینکه در i امین درج به بیش از $k = 2logn$ جستجو نیاز باشد حداکثر برابر است با:

\begin{LTR}
$2^{-2logn} = (2^{logn})^{-2} = n^{-2} = 1/n^{2}$.
\end{LTR}{}

\\
فرض کنید که A رویداد $X > 2logn $ باشد و برای $i = 1,2,...,n$  فرض کنید رویداد $A_i$ رویداد این باشد که $X_i > 2logn$. در قسمت بالا نشان دادیم که $Pr\{A_i\} \le 1/n^2$ برای هر $i$ . حال توجه کنید که 
\\
$Pr\{A\} \le Pr\{A_1\} + Pr\{A_2\} + ... + Pr\{A_n\}$
چرا که A اشتراک $A_i$ ها است. پس این عبارت حداکثر برابر است با :
\begin{LTR}
$Pr\{A\} \le n. \frac{1}{n^2}  = 1/n$.
\end{LTR}{}

\\
حال با توجه تعریف امید ریاضی داریم :

\begin{LTR}
$E[X] = \sum_{k=1}^{n} k.Pr\{X=k\} $ \\
$ = \sum_{k=1}^{[2logn]} k.Pr\{X=k\} + \sum_{k = [2logn] +1}^{n} k.Pr\{X=k\}$ \\
$ \le \sum_{k=1}^{[2logn]} [2logn].Pr\{X=k\}  + \sum_{k=[2logn]+1}^{n} n.Pr\{X = k\}$ \\
$ = [2logn] \sum_{k=1}^{[2logn]} Pr\{X=k\} +  n \sum_{k=[2logn] +1}^{n} Pr\{X = k\}.$ \\
\end{LTR}

حال برای عبارت بالا جای جمع اول عدد ۱ و جای جمع دوم با توجه به نکات ذکر شده در بالا عبارت $1/n$ را تخمین می‌زنیم داریم:
\begin{LTR}
$E[x] \le [2logn].1 + n.(1/n) $ \\
$ = [2logn] + 1$ \\
$ = O(logn) $.
\end{LTR}

}
%اگر مسئله چند قسمت داشت به این صورت عمل کنید:
%\شروع{زیرمسائل}
%\زیرمسئله[0] % مسئلۀ ۱آ
%\زیرمسئله[0] % مسئلۀ ۱ب
%\زیرمسئله[0] % مسئلۀ ۱پ
%\زیرمسئله[0] % مسئلۀ ۱ت
%\پایان{زیرمسائل}



%%%%%%%%%%%%%%%%%%%%%%%%%%%%%%%%%%%%%%%%%%%%%%%%%%%%%
% برخی دستورات مفید برای تمرین‌ها. %
%%%%%%%%%%%%%%%%%%%%%%%%%%%%%%%%%%%%%%%%%%%%%%%%%%%%%

% فرمول‌های نمونه:
%$f(x) = \Theta(x)$
%$T(x, y) \leq \log(x) + 2^y + \binom{2n}{n}$
% {\tt code\_function}


% لیست موارد بدون شماره‌گذاری:
%\شروع{فقرات}
%    \فقره مورد اول 
%        \شروع{فقرات}
%            \فقره مورد اول داخلی 
%                \شروع{فقرات}
%                    \فقره اولین مورد خیلی داخلی 
%                    \فقره دومین مورد خیلی داخلی 
%                \پایان{فقرات}
%            \فقره مورد دوم داخلی 
%        \پایان{فقرات}
%    \فقره مورد دوم 
%\پایان{فقرات}

% لیست موارد شماره‌گذاری شده:
%\شروع{شمارش}
%    \فقره یک مورد 
%    \فقره یک مورد دیگر 
%    \فقره باز هم یک مورد دیگر
%\پایان{شمارش}

% برای نوشتن یک فرمول چند خطی که یک شماره بخوره. برای حالت‌های دیگر به سایت زیر مراجعه کنید:
% https://www.overleaf.com/learn/latex/Aligning_equations_with_amsmath
%\begin{align} 
%    \begin{split}
%        (x+y)^3 &= (x+y)^2(x+y) \\
%                &= (x^2+2xy+y^2)(x+y) \\
%                &= (x^3+2x^2y+xy^2) + (x^2y+2xy^2+y^3) \\
%                &= x^3+3x^2y+3xy^2+y^3
%    \end{split}                                 
%\end{align}

% نوشتن ماتریس و آرایه دوبعدی
%\begin{align}
%    A = 
%    \begin{bmatrix}
%        A_{11} & A_{21} \\
%        A_{21} & A_{22}
%    \end{bmatrix}
%\end{align}

% اضافه کردن تصویر با شرح. برای اضافه کردن تصویر بدون شرح، فقط دستور درج‌تصویر را استفاده کنید.
%\شروع{شکل}[h!]
%\تنظیم‌ازوسط
%  \درج‌تصویر[پهنا=0.5\پهنای‌متن]{img.jpg}
%\شرح{شرح تصویر}
%\برچسب{تصویر:۱}
%\پایان{شکل}
%\پایان{زیرمسائل}

%کد زیر را ببینید.
%%\lstinputlisting{source_filename.py}
%\begin{latin}
%\begin{lstlisting}
%# Import the modules
%import sys
%import random
%
%ans = True
%
%while ans:
%question = raw_input("Ask the magic 8 ball a question: (press enter to quit) ")
%
%answers = random.randint(1,8)
%
%if question == "":
%sys.exit()
%
%elif answers == 1:
%print "It is certain"
%
%elif answers == 2:
%print "Outlook good"
%
%elif answers == 3:
%print "You may rely on it"
%
%elif answers == 4:
%print "Ask again later"
%
%elif answers == 5:
%print "Concentrate and ask again"
%
%elif answers == 6:
%print "Reply hazy, try again"
%
%elif answers == 7:
%print "My reply is no"
%
%elif answers == 8:
%print "My sources say no"
%\end{lstlisting}
%\end{latin}
